\section{Introduction}\label{sec:introduction}

Many robotic tasks are posed as reaching a goal: a quadruped that should climb onto a box is rewarded for standing on top of it, not for how it gets there. Under such a sparse reward and a long horizon, a policy started from the initial state distribution essentially never reaches the goal and therefore receives no learning signal~\citep{florensa2017reverse}. Policy gradient methods such as PPO~\citep{schulman2017proximal} perform a local search that improves the policy from the returns of the trajectories it currently generates, and a reward that none of these trajectories collects cannot guide it.

In quadrupedal locomotion, terrain curricula with dense rewards~\citep{rudin2022learning, lee2020learning}, reference motions from trajectory optimization~\citep{fuchioka2023opt}, and hand-designed curricula~\citep{atanassov2024curriculum} supply the missing signal, achieving good performance on a wide range of difficult terrains and tasks. However, each requires a demonstration or engineering specific to one task, and curiosity-driven exploration~\citep{schwarke2023curiosity} adds rewards that the task itself does not provide.

Automatic curriculum learning replaces this design effort with intermediate tasks generated during training~\citep{narvekar2020curriculum, portelas2020automatic}. Start-state and goal curricula address goal-conditioned tasks directly by shortening the task from one of its two ends: start-state curricula begin closer to the goal, with start states taken from a reference trajectory~\citep{peng2018deepmimic, resnick2018backplay} or from random walks outward from the goal~\citep{florensa2017reverse}, and goal curricula keep the start fixed and propose intermediate goals~\citep{florensa2018automatic, ren2019exploration, cho2023outcome}. Both typically expand from one end only, so the curriculum has to carry its frontier across the entire distance to the end that stays fixed. The policy, in turn, sees the target task from one side, reaching the goal from many starts or many goals from the start, but not both. In addition, the random walks of reverse curriculum generation~\citep[RC;][]{florensa2017reverse} are seeded without regard to coverage, so nothing promotes the frontier into unvisited task space or toward the initial state distribution.

Two mechanisms from sampling-based motion planning address exactly this. Rapidly-exploring random trees~\citep[RRT;][]{lavalle2001randomized} extend the vertex nearest to a uniform sample of the space, so growth is biased toward large unexplored regions (the Voronoi bias), and RRT-Connect~\citep{kuffner2000rrt} grows one tree from the start and one from the goal and extends each toward the other until they meet. BVER (Bidirectional incremental Voronoi-biased Exploration curriculum for Reinforcement learning) transfers both to curriculum generation. It grows start states outward from the goal by random walks (the backwards expansion) and goals outward from the initial state distribution in the same way (the forwards expansion), connects the two sets, and trains a single goal-conditioned policy on both intermediate starts and goals. This creates a curriculum whose proposed tasks are both at the right difficulty and promote coverage of the task space, and thus of the target task.

We evaluate BVER with sparse rewards in three domains: point-mass mazes, simulated quadrupedal box climbing, and a Tower of Hanoi manipulation experiment. In all three, PPO rarely or never finds a solution without a curriculum, and BVER converges faster than every reference-free curriculum we compare against on all but the simplest maze. On box climbing, the resulting policy holds up under start, goal, and yaw variations on which the other curricula fail. Ablations on the navigation, locomotion and manipulation experiments show that expanding from both ends outperforms either direction alone.

Our contributions are:
\begin{itemize}[leftmargin=16pt, topsep=0.0em]
    \item An expansion mechanism carrying the Voronoi bias of RRT and the connect step of RRT-Connect over to curriculum generation, growing each frontier toward unexplored task space and toward the other frontier.
    \item An automatic start-goal curriculum for goal-conditioned RL with sparse rewards and long horizons that expands bidirectionally to propose starts and goals to train one policy on both, without reference motions, shaped rewards, or a hand-designed curriculum.
    \item An evaluation against reference-free and reference-based curricula on point-mass navigation, high-dimensional quadrupedal locomotion, and manipulation with sparse rewards and long horizons, and ablations BVER's components.
\end{itemize}
In total, we evaluate on thirteen experiments: three point-mass mazes, four box climbing variants, a Tower of Hanoi manipulation, the three-route terrain, and four scanned real-world terrains.
